% Ngân hàng câu hỏi Toán 10 (KNTT) -- Chương 2, Bài 3, Dạng 08
%% Dạng: Nhận biết bất phương trình bậc nhất hai ẩn, nghiệm, điểm thuộc miền nghiệm
%% 100 câu: NB 30 -- TH 40 -- VD 20 -- VDC 10 ; mỗi mức: trắc nghiệm : đúng/sai : trả lời ngắn = 3 : 1 : 1
%% Phương án ĐẦU TIÊN trong \cauTN luôn là đáp án đúng (gói trondethi tự đảo).
%% Phần II: \y{nội dung}{d|s}. Phần III: đáp án đặt ở đối số cuối.

\cauTN[id=T10-C2-B3-D08-NB-001,muc=NB,loigiai={Bất phương trình bậc nhất hai ẩn có dạng $ax+by+c<0$ (hoặc $\le,>,\ge$) với $a$, $b$ không đồng thời bằng $0$. Chỉ có $4x-5y<10$ có dạng đó.}]{Bất phương trình nào sau đây là bất phương trình bậc nhất hai ẩn?}
  {$4x-5y<10$.}
  {$3x^2-y^2<7$.}{$x+y^2>5$.}{$2x+y^2>4$.}

\cauTN[id=T10-C2-B3-D08-NB-002,muc=NB,loigiai={Thay $(-3;-5)$ vào vế trái ta được $-10$, thoả mãn $-10\le -6$. Các cặp số còn lại không thoả mãn.}]{Cặp số nào sau đây là một nghiệm của bất phương trình $5x-y\le -6$?}
  {$(-3;-5)$.}
  {$(4;3)$.}{$(0;0)$.}{$(3;-2)$.}

\cauTN[id=T10-C2-B3-D08-NB-003,muc=NB,loigiai={Thay toạ độ điểm $(-2;3)$: $-8\ge -3$ là sai, nên điểm này không thuộc miền nghiệm.}]{Điểm nào sau đây \textbf{không} thuộc miền nghiệm của bất phương trình $x-2y\ge -3$?}
  {$(-2;3)$.}
  {$(0;-4)$.}{$(-3;-5)$.}{$(1;-3)$.}

\cauTN[id=T10-C2-B3-D08-NB-004,muc=NB,loigiai={Thay toạ độ điểm $N(5;-2)$ vào bất phương trình: $26>5$ đúng.}]{Trong mặt phẳng toạ độ $Oxy$, điểm nào sau đây thuộc miền nghiệm của bất phương trình $6x+2y>5$?}
  {$N(5;-2)$.}
  {$E(-4;5)$.}{$P(-2;-4)$.}{$A(-1;-1)$.}

\cauTN[id=T10-C2-B3-D08-NB-005,muc=NB,loigiai={Bất phương trình bậc nhất hai ẩn có dạng $ax+by+c<0$ (hoặc $\le,>,\ge$) với $a$, $b$ không đồng thời bằng $0$. Chỉ có $6x+5y\ge -10$ có dạng đó.}]{Bất phương trình nào sau đây là bất phương trình bậc nhất hai ẩn?}
  {$6x+5y\ge -10$.}
  {$\dfrac{2}{x}+y<7$.}{$x^2+2y\le 7$.}{$5x^2-y^2<8$.}

\cauTN[id=T10-C2-B3-D08-NB-006,muc=NB,loigiai={Thay $(2;-4)$ vào vế trái ta được $20$, thoả mãn $20>-10$. Các cặp số còn lại không thoả mãn.}]{Cặp số nào sau đây là một nghiệm của bất phương trình $4x-3y>-10$?}
  {$(2;-4)$.}
  {$(-2;3)$.}{$(-4;1)$.}{$(-5;5)$.}

\cauTN[id=T10-C2-B3-D08-NB-007,muc=NB,loigiai={Thay toạ độ điểm $(4;-4)$: $20\le -3$ là sai, nên điểm này không thuộc miền nghiệm.}]{Điểm nào sau đây \textbf{không} thuộc miền nghiệm của bất phương trình $6x+y\le -3$?}
  {$(4;-4)$.}
  {$(-2;5)$.}{$(-5;3)$.}{$(0;-4)$.}

\cauTN[id=T10-C2-B3-D08-NB-008,muc=NB,loigiai={Thay toạ độ điểm $A(0;5)$ vào bất phương trình: $25\ge -6$ đúng.}]{Trong mặt phẳng toạ độ $Oxy$, điểm nào sau đây thuộc miền nghiệm của bất phương trình $6x+5y\ge -6$?}
  {$A(0;5)$.}
  {$D(-1;-5)$.}{$Q(-5;2)$.}{$N(0;-5)$.}

\cauTN[id=T10-C2-B3-D08-NB-009,muc=NB,loigiai={Bất phương trình bậc nhất hai ẩn có dạng $ax+by+c<0$ (hoặc $\le,>,\ge$) với $a$, $b$ không đồng thời bằng $0$. Chỉ có $x+y<1$ có dạng đó.}]{Bất phương trình nào sau đây là bất phương trình bậc nhất hai ẩn?}
  {$x+y<1$.}
  {$2x^2-y^2<8$.}{$5x+y^2>3$.}{$x+y+z\le 51$.}

\cauTN[id=T10-C2-B3-D08-NB-010,muc=NB,loigiai={Thay $(-5;2)$ vào vế trái ta được $-7$, thoả mãn $-7<-5$. Các cặp số còn lại không thoả mãn.}]{Cặp số nào sau đây là một nghiệm của bất phương trình $x-y<-5$?}
  {$(-5;2)$.}
  {$(3;-3)$.}{$(3;0)$.}{$(5;-2)$.}

\cauTN[id=T10-C2-B3-D08-NB-011,muc=NB,loigiai={Thay toạ độ điểm $(5;3)$: $22<-8$ là sai, nên điểm này không thuộc miền nghiệm.}]{Điểm nào sau đây \textbf{không} thuộc miền nghiệm của bất phương trình $5x-y<-8$?}
  {$(5;3)$.}
  {$(-3;-5)$.}{$(-4;-4)$.}{$(-2;1)$.}

\cauTN[id=T10-C2-B3-D08-NB-012,muc=NB,loigiai={Thay toạ độ điểm $E(-1;-4)$ vào bất phương trình: $-12\le -11$ đúng.}]{Trong mặt phẳng toạ độ $Oxy$, điểm nào sau đây thuộc miền nghiệm của bất phương trình $4x+2y\le -11$?}
  {$E(-1;-4)$.}
  {$B(4;3)$.}{$D(-2;3)$.}{$C(1;-2)$.}

\cauTN[id=T10-C2-B3-D08-NB-013,muc=NB,loigiai={Bất phương trình bậc nhất hai ẩn có dạng $ax+by+c<0$ (hoặc $\le,>,\ge$) với $a$, $b$ không đồng thời bằng $0$. Chỉ có $5x+y<5$ có dạng đó.}]{Bất phương trình nào sau đây là bất phương trình bậc nhất hai ẩn?}
  {$5x+y<5$.}
  {$2x+y^2>5$.}{$\dfrac{3}{x}+y<3$.}{$x+y+z\le 36$.}

\cauTN[id=T10-C2-B3-D08-NB-014,muc=NB,loigiai={Thay $(-4;0)$ vào vế trái ta được $-4$, thoả mãn $-4\le 2$. Các cặp số còn lại không thoả mãn.}]{Cặp số nào sau đây là một nghiệm của bất phương trình $x+2y\le 2$?}
  {$(-4;0)$.}
  {$(4;0)$.}{$(-3;5)$.}{$(-2;4)$.}

\cauTN[id=T10-C2-B3-D08-NB-015,muc=NB,loigiai={Thay toạ độ điểm $(3;3)$: $12<-6$ là sai, nên điểm này không thuộc miền nghiệm.}]{Điểm nào sau đây \textbf{không} thuộc miền nghiệm của bất phương trình $5x-y<-6$?}
  {$(3;3)$.}
  {$(-2;0)$.}{$(-4;3)$.}{$(-2;-3)$.}

\cauTN[id=T10-C2-B3-D08-NB-016,muc=NB,loigiai={Thay toạ độ điểm $M(1;3)$ vào bất phương trình: $-1\ge -4$ đúng.}]{Trong mặt phẳng toạ độ $Oxy$, điểm nào sau đây thuộc miền nghiệm của bất phương trình $2x-y\ge -4$?}
  {$M(1;3)$.}
  {$E(0;5)$.}{$A(-4;2)$.}{$Q(-2;2)$.}

\cauTN[id=T10-C2-B3-D08-NB-017,muc=NB,loigiai={Bất phương trình bậc nhất hai ẩn có dạng $ax+by+c<0$ (hoặc $\le,>,\ge$) với $a$, $b$ không đồng thời bằng $0$. Chỉ có $4x-y>-12$ có dạng đó.}]{Bất phương trình nào sau đây là bất phương trình bậc nhất hai ẩn?}
  {$4x-y>-12$.}
  {$xy-2x\ge 9$.}{$xy-3x\ge 4$.}{$x^2+4y\le 3$.}

\cauTN[id=T10-C2-B3-D08-NB-018,muc=NB,loigiai={Thay $(0;2)$ vào vế trái ta được $-10$, thoả mãn $-10<-2$. Các cặp số còn lại không thoả mãn.}]{Cặp số nào sau đây là một nghiệm của bất phương trình $5x-5y<-2$?}
  {$(0;2)$.}
  {$(0;-1)$.}{$(5;4)$.}{$(0;-3)$.}

\cauDS[id=T10-C2-B3-D08-NB-019,muc=NB,loigiai={Với $(5;1)$: $8<7$ sai. Với $(2;-1)$: $-1<7$ đúng. Với $(4;5)$: $19<7$ sai. Với $(2;1)$: $5<7$ đúng. a) Sai. b) Đúng. c) Sai. d) Đúng.}]{Cho bất phương trình bậc nhất hai ẩn $x+3y<7$.}{%
  \y{Cặp số $(5;1)$ là một nghiệm của bất phương trình.}{s}
  \y{Cặp số $(2;-1)$ là một nghiệm của bất phương trình.}{d}
  \y{Cặp số $(4;5)$ là một nghiệm của bất phương trình.}{s}
  \y{Cặp số $(2;1)$ là một nghiệm của bất phương trình.}{d}}

\cauDS[id=T10-C2-B3-D08-NB-020,muc=NB,loigiai={Với $(1;-4)$: $-13\le -3$ đúng. Với $(2;2)$: $14\le -3$ sai. Với $(-5;-3)$: $-27\le -3$ đúng. Với $(3;1)$: $13\le -3$ sai. a) Đúng. b) Sai. c) Đúng. d) Sai.}]{Cho bất phương trình bậc nhất hai ẩn $3x+4y\le -3$.}{%
  \y{Cặp số $(1;-4)$ là một nghiệm của bất phương trình.}{d}
  \y{Cặp số $(2;2)$ là một nghiệm của bất phương trình.}{s}
  \y{Cặp số $(-5;-3)$ là một nghiệm của bất phương trình.}{d}
  \y{Cặp số $(3;1)$ là một nghiệm của bất phương trình.}{s}}

\cauDS[id=T10-C2-B3-D08-NB-021,muc=NB,loigiai={Với $(-3;-1)$: $-1\ge -6$ đúng. Với $(-1;3)$: $-17\ge -6$ sai. Với $(1;-1)$: $7\ge -6$ đúng. Với $(1;3)$: $-13\ge -6$ sai. a) Đúng. b) Sai. c) Đúng. d) Sai.}]{Cho bất phương trình bậc nhất hai ẩn $2x-5y\ge -6$.}{%
  \y{Cặp số $(-3;-1)$ là một nghiệm của bất phương trình.}{d}
  \y{Cặp số $(-1;3)$ là một nghiệm của bất phương trình.}{s}
  \y{Cặp số $(1;-1)$ là một nghiệm của bất phương trình.}{d}
  \y{Cặp số $(1;3)$ là một nghiệm của bất phương trình.}{s}}

\cauDS[id=T10-C2-B3-D08-NB-022,muc=NB,loigiai={Với $(-3;3)$: $-15>-1$ sai. Với $(-1;5)$: $-17>-1$ sai. Với $(5;-2)$: $16>-1$ đúng. Với $(0;-5)$: $15>-1$ đúng. a) Sai. b) Sai. c) Đúng. d) Đúng.}]{Cho bất phương trình bậc nhất hai ẩn $2x-3y>-1$.}{%
  \y{Cặp số $(-3;3)$ là một nghiệm của bất phương trình.}{s}
  \y{Cặp số $(-1;5)$ là một nghiệm của bất phương trình.}{s}
  \y{Cặp số $(5;-2)$ là một nghiệm của bất phương trình.}{d}
  \y{Cặp số $(0;-5)$ là một nghiệm của bất phương trình.}{d}}

\cauDS[id=T10-C2-B3-D08-NB-023,muc=NB,loigiai={Với $(-2;5)$: $-9\ge 1$ sai. Với $(4;-3)$: $11\ge 1$ đúng. Với $(-4;2)$: $-10\ge 1$ sai. Với $(2;-1)$: $5\ge 1$ đúng. a) Sai. b) Đúng. c) Sai. d) Đúng.}]{Cho bất phương trình bậc nhất hai ẩn $2x-y\ge 1$.}{%
  \y{Cặp số $(-2;5)$ là một nghiệm của bất phương trình.}{s}
  \y{Cặp số $(4;-3)$ là một nghiệm của bất phương trình.}{d}
  \y{Cặp số $(-4;2)$ là một nghiệm của bất phương trình.}{s}
  \y{Cặp số $(2;-1)$ là một nghiệm của bất phương trình.}{d}}

\cauDS[id=T10-C2-B3-D08-NB-024,muc=NB,loigiai={Với $(5;3)$: $11\le 9$ sai. Với $(5;4)$: $13\le 9$ sai. Với $(4;1)$: $6\le 9$ đúng. Với $(-5;-5)$: $-15\le 9$ đúng. a) Sai. b) Sai. c) Đúng. d) Đúng.}]{Cho bất phương trình bậc nhất hai ẩn $x+2y\le 9$.}{%
  \y{Cặp số $(5;3)$ là một nghiệm của bất phương trình.}{s}
  \y{Cặp số $(5;4)$ là một nghiệm của bất phương trình.}{s}
  \y{Cặp số $(4;1)$ là một nghiệm của bất phương trình.}{d}
  \y{Cặp số $(-5;-5)$ là một nghiệm của bất phương trình.}{d}}

\cauTLN[id=T10-C2-B3-D08-NB-025,muc=NB,loigiai={Thay lần lượt toạ độ các điểm; các điểm thoả mãn là $C$, $D$, $E$.}]{Cho bất phương trình $6x+y\ge -12$ và các điểm $A(-4;-2)$, $B(-5;-3)$, $C(4;3)$, $D(1;-1)$, $E(0;-2)$. Có bao nhiêu điểm trong các điểm trên thuộc miền nghiệm của bất phương trình?}{3}

\cauTLN[id=T10-C2-B3-D08-NB-026,muc=NB,loigiai={Thay lần lượt toạ độ các điểm; các điểm thoả mãn là $A$, $B$, $C$, $D$.}]{Cho bất phương trình $x+3y\ge 6$ và các điểm $A(4;3)$, $B(-4;4)$, $C(4;2)$, $D(4;5)$, $E(-2;-1)$. Có bao nhiêu điểm trong các điểm trên thuộc miền nghiệm của bất phương trình?}{4}

\cauTLN[id=T10-C2-B3-D08-NB-027,muc=NB,loigiai={Thay lần lượt toạ độ các điểm; các điểm thoả mãn là $A$, $B$, $C$, $D$.}]{Cho bất phương trình $5x+4y\ge -7$ và các điểm $A(3;-2)$, $B(2;5)$, $C(-2;2)$, $D(-5;5)$, $E(-4;-4)$. Có bao nhiêu điểm trong các điểm trên thuộc miền nghiệm của bất phương trình?}{4}

\cauTLN[id=T10-C2-B3-D08-NB-028,muc=NB,loigiai={Thay lần lượt toạ độ các điểm; các điểm thoả mãn là $A$, $D$, $E$.}]{Cho bất phương trình $2x-4y\ge -3$ và các điểm $A(3;-1)$, $B(-2;4)$, $C(4;3)$, $D(3;-3)$, $E(-1;-2)$. Có bao nhiêu điểm trong các điểm trên thuộc miền nghiệm của bất phương trình?}{3}

\cauTLN[id=T10-C2-B3-D08-NB-029,muc=NB,loigiai={Thay lần lượt toạ độ các điểm; các điểm thoả mãn là $A$, $B$, $C$, $E$.}]{Cho bất phương trình $3x-5y\le -5$ và các điểm $A(2;5)$, $B(-1;5)$, $C(2;3)$, $D(0;-2)$, $E(-5;2)$. Có bao nhiêu điểm trong các điểm trên thuộc miền nghiệm của bất phương trình?}{4}

\cauTLN[id=T10-C2-B3-D08-NB-030,muc=NB,loigiai={Thay lần lượt toạ độ các điểm; các điểm thoả mãn là $A$, $C$, $D$, $E$.}]{Cho bất phương trình $x+5y<5$ và các điểm $A(-4;-5)$, $B(4;5)$, $C(1;-4)$, $D(-2;-3)$, $E(2;-4)$. Có bao nhiêu điểm trong các điểm trên thuộc miền nghiệm của bất phương trình?}{4}

\cauTN[id=T10-C2-B3-D08-TH-031,muc=TH,loigiai={Số tiền phải trả là $5x+6y$ (nghìn đồng), không vượt quá $100$ nên $5x+6y\le 100$.}]{Bạn Cường có không quá $100$ nghìn đồng để mua đồ dùng. Giá mỗi quyển vở là $5$ nghìn đồng, giá mỗi cây bút là $6$ nghìn đồng. Gọi $x$, $y$ lần lượt là số quyển vở và số cây bút bạn Cường mua. Bất phương trình mô tả điều kiện về số tiền là}
  {$5x+6y\le 100$.}
  {$5x+6y\ge 100$.}{$6x+5y\le 100$.}{$5x+6y<100$.}

\cauTN[id=T10-C2-B3-D08-TH-032,muc=TH,loigiai={Thay $y=-2$: $x-4\le 3$, giải ra và lấy số nguyên lớn nhất là $x=7$.}]{Cho bất phương trình $x+2y\le 3$. Số nguyên $x$ lớn nhất sao cho cặp số $(x;-2)$ là nghiệm của bất phương trình là}
  {$7$.}
  {$11$.}{$5$.}{$8$.}

\cauTN[id=T10-C2-B3-D08-TH-033,muc=TH,loigiai={Thay $(-4;1)$ vào vế trái ta được $-13$, thoả mãn $-13\le -12$. Các cặp số còn lại không thoả mãn.}]{Cặp số nào sau đây là một nghiệm của bất phương trình $4x+3y\le -12$?}
  {$(-4;1)$.}
  {$(-2;5)$.}{$(2;-3)$.}{$(2;-4)$.}

\cauTN[id=T10-C2-B3-D08-TH-034,muc=TH,loigiai={Chuyển vế, rút gọn: $3x+5y\ge -13$.}]{Bất phương trình $5(x+2)+5y\ge 2x-3$ tương đương với bất phương trình nào sau đây?}
  {$3x+5y\ge -13$.}
  {$3x+5y\le -13$.}{$3x+5y\ge 13$.}{$3x-5y\ge -13$.}

\cauTN[id=T10-C2-B3-D08-TH-035,muc=TH,loigiai={Số tiền phải trả là $10x+15y$ (nghìn đồng), không vượt quá $180$ nên $10x+15y\le 180$.}]{Bạn Hà có không quá $180$ nghìn đồng để mua đồ dùng. Giá mỗi quyển vở là $10$ nghìn đồng, giá mỗi cây bút là $15$ nghìn đồng. Gọi $x$, $y$ lần lượt là số quyển vở và số cây bút bạn Hà mua. Bất phương trình mô tả điều kiện về số tiền là}
  {$10x+15y\le 180$.}
  {$10x+15y\ge 180$.}{$15x+10y\le 180$.}{$10x+15y<180$.}

\cauTN[id=T10-C2-B3-D08-TH-036,muc=TH,loigiai={Thay $y=3$: $4x+15\le -4$, giải ra và lấy số nguyên lớn nhất là $x=-5$.}]{Cho bất phương trình $4x+5y\le -4$. Số nguyên $x$ lớn nhất sao cho cặp số $(x;3)$ là nghiệm của bất phương trình là}
  {$-5$.}
  {$-8$.}{$-6$.}{$-3$.}

\cauTN[id=T10-C2-B3-D08-TH-037,muc=TH,loigiai={Thay $(5;-5)$ vào vế trái ta được $-5$, thoả mãn $-5>-7$. Các cặp số còn lại không thoả mãn.}]{Cặp số nào sau đây là một nghiệm của bất phương trình $3x+4y>-7$?}
  {$(5;-5)$.}
  {$(-5;0)$.}{$(-5;-5)$.}{$(-2;-3)$.}

\cauTN[id=T10-C2-B3-D08-TH-038,muc=TH,loigiai={Chuyển vế, rút gọn: $x-2y\le 3$.}]{Bất phương trình $3x+2y\ge 4x-3$ tương đương với bất phương trình nào sau đây?}
  {$x-2y\le 3$.}
  {$x-2y\ge 3$.}{$x-2y\le -3$.}{$x+2y\le 3$.}

\cauTN[id=T10-C2-B3-D08-TH-039,muc=TH,loigiai={Số tiền phải trả là $15x+10y$ (nghìn đồng), không vượt quá $260$ nên $15x+10y\le 260$.}]{Bạn Mai có không quá $260$ nghìn đồng để mua đồ dùng. Giá mỗi cuốn truyện là $15$ nghìn đồng, giá mỗi tấm bưu thiếp là $10$ nghìn đồng. Gọi $x$, $y$ lần lượt là số cuốn truyện và số tấm bưu thiếp bạn Mai mua. Bất phương trình mô tả điều kiện về số tiền là}
  {$15x+10y\le 260$.}
  {$15x+10y\ge 260$.}{$10x+15y\le 260$.}{$15x+10y<260$.}

\cauTN[id=T10-C2-B3-D08-TH-040,muc=TH,loigiai={Thay $y=3$: $4x+12>11$, giải ra và lấy số nguyên nhỏ nhất là $x=0$.}]{Cho bất phương trình $4x+4y>11$. Số nguyên $x$ nhỏ nhất sao cho cặp số $(x;3)$ là nghiệm của bất phương trình là}
  {$0$.}
  {$1$.}{$-3$.}{$2$.}

\cauTN[id=T10-C2-B3-D08-TH-041,muc=TH,loigiai={Thay $(2;-3)$ vào vế trái ta được $3$, thoả mãn $3\le 7$. Các cặp số còn lại không thoả mãn.}]{Cặp số nào sau đây là một nghiệm của bất phương trình $3x+y\le 7$?}
  {$(2;-3)$.}
  {$(3;1)$.}{$(4;-1)$.}{$(3;5)$.}

\cauTN[id=T10-C2-B3-D08-TH-042,muc=TH,loigiai={Số tiền phải trả là $6x+20y$ (nghìn đồng), không vượt quá $160$ nên $6x+20y\le 160$.}]{Bạn Hải có không quá $160$ nghìn đồng để mua đồ dùng. Giá mỗi kg táo là $6$ nghìn đồng, giá mỗi kg cam là $20$ nghìn đồng. Gọi $x$, $y$ lần lượt là số kg táo và số kg cam bạn Hải mua. Bất phương trình mô tả điều kiện về số tiền là}
  {$6x+20y\le 160$.}
  {$6x+20y\ge 160$.}{$20x+6y\le 160$.}{$6x+20y<160$.}

\cauTN[id=T10-C2-B3-D08-TH-043,muc=TH,loigiai={Thay $y=-1$: $3x+5\ge 9$, giải ra và lấy số nguyên nhỏ nhất là $x=2$.}]{Cho bất phương trình $3x-5y\ge 9$. Số nguyên $x$ nhỏ nhất sao cho cặp số $(x;-1)$ là nghiệm của bất phương trình là}
  {$2$.}
  {$3$.}{$1$.}{$6$.}

\cauTN[id=T10-C2-B3-D08-TH-044,muc=TH,loigiai={Thay $(0;1)$ vào vế trái ta được $-6$, thoả mãn $-6\ge -12$. Các cặp số còn lại không thoả mãn.}]{Cặp số nào sau đây là một nghiệm của bất phương trình $x-6y\ge -12$?}
  {$(0;1)$.}
  {$(2;4)$.}{$(-4;2)$.}{$(5;5)$.}

\cauTN[id=T10-C2-B3-D08-TH-045,muc=TH,loigiai={Chuyển vế, rút gọn: $x+y\le -11$.}]{Bất phương trình $3(x+3)+y\le 2x-2$ tương đương với bất phương trình nào sau đây?}
  {$x+y\le -11$.}
  {$x+y\ge -11$.}{$x+y\le 11$.}{$x-y\le -11$.}

\cauTN[id=T10-C2-B3-D08-TH-046,muc=TH,loigiai={Số tiền phải trả là $20x+10y$ (nghìn đồng), không vượt quá $240$ nên $20x+10y\le 240$.}]{Bạn Minh có không quá $240$ nghìn đồng để mua đồ dùng. Giá mỗi quyển vở là $20$ nghìn đồng, giá mỗi cây bút là $10$ nghìn đồng. Gọi $x$, $y$ lần lượt là số quyển vở và số cây bút bạn Minh mua. Bất phương trình mô tả điều kiện về số tiền là}
  {$20x+10y\le 240$.}
  {$20x+10y\ge 240$.}{$10x+20y\le 240$.}{$20x+10y<240$.}

\cauTN[id=T10-C2-B3-D08-TH-047,muc=TH,loigiai={Thay $y=-2$: $3x-4<-2$, giải ra và lấy số nguyên lớn nhất là $x=0$.}]{Cho bất phương trình $3x+2y<-2$. Số nguyên $x$ lớn nhất sao cho cặp số $(x;-2)$ là nghiệm của bất phương trình là}
  {$0$.}
  {$-2$.}{$4$.}{$3$.}

\cauTN[id=T10-C2-B3-D08-TH-048,muc=TH,loigiai={Thay $(-2;1)$ vào vế trái ta được $-13$, thoả mãn $-13<4$. Các cặp số còn lại không thoả mãn.}]{Cặp số nào sau đây là một nghiệm của bất phương trình $5x-3y<4$?}
  {$(-2;1)$.}
  {$(3;3)$.}{$(3;-4)$.}{$(5;-4)$.}

\cauTN[id=T10-C2-B3-D08-TH-049,muc=TH,loigiai={Chuyển vế, rút gọn: $2x-y<-13$.}]{Bất phương trình $3(x-5)+y>5x-2$ tương đương với bất phương trình nào sau đây?}
  {$2x-y<-13$.}
  {$2x-y>-13$.}{$2x-y<13$.}{$2x+y<-13$.}

\cauTN[id=T10-C2-B3-D08-TH-050,muc=TH,loigiai={Số tiền phải trả là $12x+15y$ (nghìn đồng), không vượt quá $140$ nên $12x+15y\le 140$.}]{Bạn Nam có không quá $140$ nghìn đồng để mua đồ dùng. Giá mỗi cuốn truyện là $12$ nghìn đồng, giá mỗi tấm bưu thiếp là $15$ nghìn đồng. Gọi $x$, $y$ lần lượt là số cuốn truyện và số tấm bưu thiếp bạn Nam mua. Bất phương trình mô tả điều kiện về số tiền là}
  {$12x+15y\le 140$.}
  {$12x+15y\ge 140$.}{$15x+12y\le 140$.}{$12x+15y<140$.}

\cauTN[id=T10-C2-B3-D08-TH-051,muc=TH,loigiai={Thay $y=0$: $4x\le 11$, giải ra và lấy số nguyên lớn nhất là $x=2$.}]{Cho bất phương trình $4x-y\le 11$. Số nguyên $x$ lớn nhất sao cho cặp số $(x;0)$ là nghiệm của bất phương trình là}
  {$2$.}
  {$5$.}{$0$.}{$1$.}

\cauTN[id=T10-C2-B3-D08-TH-052,muc=TH,loigiai={Thay $(5;0)$ vào vế trái ta được $25$, thoả mãn $25>-1$. Các cặp số còn lại không thoả mãn.}]{Cặp số nào sau đây là một nghiệm của bất phương trình $5x-3y>-1$?}
  {$(5;0)$.}
  {$(-1;1)$.}{$(-3;3)$.}{$(-5;5)$.}

\cauTN[id=T10-C2-B3-D08-TH-053,muc=TH,loigiai={Chuyển vế, rút gọn: $3x-3y\ge -2$.}]{Bất phương trình $(x-4)+3y\le 4x-2$ tương đương với bất phương trình nào sau đây?}
  {$3x-3y\ge -2$.}
  {$3x-3y\le -2$.}{$3x-3y\ge 2$.}{$3x+3y\ge -2$.}

\cauTN[id=T10-C2-B3-D08-TH-054,muc=TH,loigiai={Số tiền phải trả là $6x+12y$ (nghìn đồng), không vượt quá $260$ nên $6x+12y\le 260$.}]{Bạn Sơn có không quá $260$ nghìn đồng để mua đồ dùng. Giá mỗi kg táo là $6$ nghìn đồng, giá mỗi kg cam là $12$ nghìn đồng. Gọi $x$, $y$ lần lượt là số kg táo và số kg cam bạn Sơn mua. Bất phương trình mô tả điều kiện về số tiền là}
  {$6x+12y\le 260$.}
  {$6x+12y\ge 260$.}{$12x+6y\le 260$.}{$6x+12y<260$.}

\cauDS[id=T10-C2-B3-D08-TH-055,muc=TH,loigiai={Doanh thu $4x+6y\ge 110$. Với $(6;3)$: $42$; với $(14;4)$: $80$. Chỉ bán trà đào: $6y\ge 110\Leftrightarrow y\ge \dfrac{55}{3}$, nên cần ít nhất $19$ cốc. a) Đúng. b) Sai. c) Sai. d) Sai.}]{Một nhóm học sinh tổ chức gian hàng ở hội chợ xuân của trường, bán hai loại nước: nước chanh giá $4$ nghìn đồng/cốc và trà đào giá $6$ nghìn đồng/cốc. Nhóm đặt mục tiêu doanh thu ít nhất $110$ nghìn đồng. Gọi $x$, $y$ lần lượt là số cốc nước chanh và trà đào bán được.}{%
  \y{Điều kiện về doanh thu là $4x+6y\ge 110$.}{d}
  \y{Nếu bán được $6$ cốc nước chanh và $3$ cốc trà đào thì nhóm đạt mục tiêu.}{s}
  \y{Nếu bán được $14$ cốc nước chanh và $4$ cốc trà đào thì nhóm đạt mục tiêu.}{s}
  \y{Nếu chỉ bán trà đào thì cần bán ít nhất $20$ cốc để đạt mục tiêu.}{s}}

\cauDS[id=T10-C2-B3-D08-TH-056,muc=TH,loigiai={Với $(-1;-4)$: $7\ge -1$ đúng. Với $(-2;-2)$: $-4\ge -1$ sai. Với $(-3;5)$: $-30\ge -1$ sai. Với $(1;0)$: $5\ge -1$ đúng. a) Đúng. b) Sai. c) Sai. d) Đúng.}]{Cho bất phương trình bậc nhất hai ẩn $5x-3y\ge -1$.}{%
  \y{Cặp số $(-1;-4)$ là một nghiệm của bất phương trình.}{d}
  \y{Cặp số $(-2;-2)$ là một nghiệm của bất phương trình.}{s}
  \y{Cặp số $(-3;5)$ là một nghiệm của bất phương trình.}{s}
  \y{Cặp số $(1;0)$ là một nghiệm của bất phương trình.}{d}}

\cauDS[id=T10-C2-B3-D08-TH-057,muc=TH,loigiai={Doanh thu $6x+8y\ge 60$. Với $(2;13)$: $116$; với $(1;3)$: $30$. Chỉ bán trà đào: $8y\ge 60\Leftrightarrow y\ge \dfrac{15}{2}$, nên cần ít nhất $8$ cốc. a) Đúng. b) Đúng. c) Sai. d) Sai.}]{Một nhóm học sinh tổ chức gian hàng ở hội chợ xuân của trường, bán hai loại nước: nước chanh giá $6$ nghìn đồng/cốc và trà đào giá $8$ nghìn đồng/cốc. Nhóm đặt mục tiêu doanh thu ít nhất $60$ nghìn đồng. Gọi $x$, $y$ lần lượt là số cốc nước chanh và trà đào bán được.}{%
  \y{Điều kiện về doanh thu là $6x+8y\ge 60$.}{d}
  \y{Nếu bán được $2$ cốc nước chanh và $13$ cốc trà đào thì nhóm đạt mục tiêu.}{d}
  \y{Nếu bán được $1$ cốc nước chanh và $3$ cốc trà đào thì nhóm đạt mục tiêu.}{s}
  \y{Nếu chỉ bán trà đào thì cần bán ít nhất $9$ cốc để đạt mục tiêu.}{s}}

\cauDS[id=T10-C2-B3-D08-TH-058,muc=TH,loigiai={Với $(4;-4)$: $20<-1$ sai. Với $(1;3)$: $9<-1$ sai. Với $(-1;-5)$: $-11<-1$ đúng. Với $(-1;1)$: $-5<-1$ đúng. a) Sai. b) Sai. c) Đúng. d) Đúng.}]{Cho bất phương trình bậc nhất hai ẩn $6x+y<-1$.}{%
  \y{Cặp số $(4;-4)$ là một nghiệm của bất phương trình.}{s}
  \y{Cặp số $(1;3)$ là một nghiệm của bất phương trình.}{s}
  \y{Cặp số $(-1;-5)$ là một nghiệm của bất phương trình.}{d}
  \y{Cặp số $(-1;1)$ là một nghiệm của bất phương trình.}{d}}

\cauDS[id=T10-C2-B3-D08-TH-059,muc=TH,loigiai={Doanh thu $4x+6y\ge 90$. Với $(0;11)$: $66$; với $(0;12)$: $72$. Chỉ bán trà đào: $6y\ge 90\Leftrightarrow y\ge 15$, nên cần ít nhất $15$ cốc. a) Đúng. b) Sai. c) Sai. d) Đúng.}]{Một nhóm học sinh tổ chức gian hàng ở hội chợ xuân của trường, bán hai loại nước: nước chanh giá $4$ nghìn đồng/cốc và trà đào giá $6$ nghìn đồng/cốc. Nhóm đặt mục tiêu doanh thu ít nhất $90$ nghìn đồng. Gọi $x$, $y$ lần lượt là số cốc nước chanh và trà đào bán được.}{%
  \y{Điều kiện về doanh thu là $4x+6y\ge 90$.}{d}
  \y{Nếu bán được $0$ cốc nước chanh và $11$ cốc trà đào thì nhóm đạt mục tiêu.}{s}
  \y{Nếu bán được $0$ cốc nước chanh và $12$ cốc trà đào thì nhóm đạt mục tiêu.}{s}
  \y{Nếu chỉ bán trà đào thì cần bán ít nhất $15$ cốc để đạt mục tiêu.}{d}}

\cauDS[id=T10-C2-B3-D08-TH-060,muc=TH,loigiai={Với $(2;2)$: $16\ge 5$ đúng. Với $(-2;0)$: $-10\ge 5$ sai. Với $(5;2)$: $31\ge 5$ đúng. Với $(-5;-4)$: $-37\ge 5$ sai. a) Đúng. b) Sai. c) Đúng. d) Sai.}]{Cho bất phương trình bậc nhất hai ẩn $5x+3y\ge 5$.}{%
  \y{Cặp số $(2;2)$ là một nghiệm của bất phương trình.}{d}
  \y{Cặp số $(-2;0)$ là một nghiệm của bất phương trình.}{s}
  \y{Cặp số $(5;2)$ là một nghiệm của bất phương trình.}{d}
  \y{Cặp số $(-5;-4)$ là một nghiệm của bất phương trình.}{s}}

\cauDS[id=T10-C2-B3-D08-TH-061,muc=TH,loigiai={Doanh thu $12x+6y\ge 60$. Với $(11;11)$: $198$; với $(11;1)$: $138$. Chỉ bán trà đào: $6y\ge 60\Leftrightarrow y\ge 10$, nên cần ít nhất $10$ cốc. a) Đúng. b) Đúng. c) Đúng. d) Sai.}]{Một nhóm học sinh tổ chức gian hàng ở hội chợ xuân của trường, bán hai loại nước: nước chanh giá $12$ nghìn đồng/cốc và trà đào giá $6$ nghìn đồng/cốc. Nhóm đặt mục tiêu doanh thu ít nhất $60$ nghìn đồng. Gọi $x$, $y$ lần lượt là số cốc nước chanh và trà đào bán được.}{%
  \y{Điều kiện về doanh thu là $12x+6y\ge 60$.}{d}
  \y{Nếu bán được $11$ cốc nước chanh và $11$ cốc trà đào thì nhóm đạt mục tiêu.}{d}
  \y{Nếu bán được $11$ cốc nước chanh và $1$ cốc trà đào thì nhóm đạt mục tiêu.}{d}
  \y{Nếu chỉ bán trà đào thì cần bán ít nhất $11$ cốc để đạt mục tiêu.}{s}}

\cauDS[id=T10-C2-B3-D08-TH-062,muc=TH,loigiai={Với $(2;5)$: $25\le -1$ sai. Với $(-4;-1)$: $-23\le -1$ đúng. Với $(0;-5)$: $-15\le -1$ đúng. Với $(5;4)$: $37\le -1$ sai. a) Sai. b) Đúng. c) Đúng. d) Sai.}]{Cho bất phương trình bậc nhất hai ẩn $5x+3y\le -1$.}{%
  \y{Cặp số $(2;5)$ là một nghiệm của bất phương trình.}{s}
  \y{Cặp số $(-4;-1)$ là một nghiệm của bất phương trình.}{d}
  \y{Cặp số $(0;-5)$ là một nghiệm của bất phương trình.}{d}
  \y{Cặp số $(5;4)$ là một nghiệm của bất phương trình.}{s}}

\cauTLN[id=T10-C2-B3-D08-TH-063,muc=TH,loigiai={Xét lần lượt $x=0,1,\ldots,1$ và đếm số giá trị $y$ tự nhiên tương ứng: 2+1=3.}]{Có bao nhiêu cặp số tự nhiên $(x;y)$ thoả mãn bất phương trình $3x+4y\le 4$?}{3}

\cauTLN[id=T10-C2-B3-D08-TH-064,muc=TH,loigiai={Thay $y=0$ được bất phương trình một ẩn $x$; các số nguyên thoả mãn trong đoạn $[-10;10]$ có $11$ số.}]{Cho bất phương trình $3x+y>-2$. Có bao nhiêu số nguyên $x\in[-10;10]$ để cặp số $(x;0)$ là nghiệm của bất phương trình?}{11}

\cauTLN[id=T10-C2-B3-D08-TH-065,muc=TH,loigiai={Xét lần lượt $x=0,1,\ldots,1$ và đếm số giá trị $y$ tự nhiên tương ứng: 3+1=4.}]{Có bao nhiêu cặp số tự nhiên $(x;y)$ thoả mãn bất phương trình $3x+2y\le 4$?}{4}

\cauTLN[id=T10-C2-B3-D08-TH-066,muc=TH,loigiai={Thay $y=1$ được bất phương trình một ẩn $x$; các số nguyên thoả mãn trong đoạn $[-10;10]$ có $10$ số.}]{Cho bất phương trình $4x-y>0$. Có bao nhiêu số nguyên $x\in[-10;10]$ để cặp số $(x;1)$ là nghiệm của bất phương trình?}{10}

\cauTLN[id=T10-C2-B3-D08-TH-067,muc=TH,loigiai={Xét lần lượt $x=0,1,\ldots,9$ và đếm số giá trị $y$ tự nhiên tương ứng: 3+3+2+2+2+2+1+1+1+1=18.}]{Có bao nhiêu cặp số tự nhiên $(x;y)$ thoả mãn bất phương trình $x+4y\le 9$?}{18}

\cauTLN[id=T10-C2-B3-D08-TH-068,muc=TH,loigiai={Thay $y=2$ được bất phương trình một ẩn $x$; các số nguyên thoả mãn trong đoạn $[-10;10]$ có $17$ số.}]{Cho bất phương trình $x-y<5$. Có bao nhiêu số nguyên $x\in[-10;10]$ để cặp số $(x;2)$ là nghiệm của bất phương trình?}{17}

\cauTLN[id=T10-C2-B3-D08-TH-069,muc=TH,loigiai={Xét lần lượt $x=0,1,\ldots,3$ và đếm số giá trị $y$ tự nhiên tương ứng: 3+2+2+1=8.}]{Có bao nhiêu cặp số tự nhiên $(x;y)$ thoả mãn bất phương trình $3x+4y\le 10$?}{8}

\cauTLN[id=T10-C2-B3-D08-TH-070,muc=TH,loigiai={Thay $y=-3$ được bất phương trình một ẩn $x$; các số nguyên thoả mãn trong đoạn $[-10;10]$ có $13$ số.}]{Cho bất phương trình $x-3y>6$. Có bao nhiêu số nguyên $x\in[-10;10]$ để cặp số $(x;-3)$ là nghiệm của bất phương trình?}{13}

\cauTN[id=T10-C2-B3-D08-VD-071,muc=VD,loigiai={Điều kiện: $15x+45y\le 390$. Với $x=3$: $45y\le 345$, suy ra $y\le \dfrac{23}{3}$. Vậy nhiều nhất $7$ buổi.}]{Trong một tuần, bạn Uyên dành thời gian tự học hai môn Toán và Tiếng Anh. Mỗi buổi học Toán kéo dài $15$ phút, mỗi buổi học Tiếng Anh kéo dài $45$ phút. Vì còn phải tham gia câu lạc bộ thể thao, tổng thời gian tự học hai môn này trong tuần của Uyên không vượt quá $390$ phút. Gọi $x$, $y$ lần lượt là số buổi học Toán và Tiếng Anh trong tuần. Nếu Uyên đã lên lịch $3$ buổi học Toán thì số buổi học Tiếng Anh nhiều nhất có thể là}
  {$7$.}
  {$23$.}{$9$.}{$8$.}

\cauTN[id=T10-C2-B3-D08-VD-072,muc=VD,loigiai={Điều kiện $90x+180y\ge 900$. Với phương án $(2;2)$: $540<900$ nên không đáp ứng.}]{Một bạn học sinh lớp 10 muốn bổ sung năng lượng cho bữa sáng trước giờ học thể dục. Mỗi hộp sữa chua cung cấp khoảng $90$ kcal, mỗi lát bánh mì nguyên cám cung cấp khoảng $180$ kcal. Theo tư vấn của nhân viên y tế trường, bữa sáng của bạn cần cung cấp không ít hơn $900$ kcal. Gọi $x$ là số hộp sữa chua và $y$ là số lát bánh mì bạn ăn. Phương án nào sau đây \textbf{không} đáp ứng yêu cầu về năng lượng?}
  {$2$ hộp sữa chua và $2$ lát bánh mì.}
  {$4$ hộp sữa chua và $3$ lát bánh mì.}{$2$ hộp sữa chua và $5$ lát bánh mì.}{$3$ hộp sữa chua và $4$ lát bánh mì.}

\cauTN[id=T10-C2-B3-D08-VD-073,muc=VD,loigiai={Điều kiện: $20x+45y\le 630$. Với $x=3$: $45y\le 570$, suy ra $y\le \dfrac{38}{3}$. Vậy nhiều nhất $12$ buổi.}]{Trong một tuần, bạn Nhung dành thời gian tự học hai môn Toán và Tiếng Anh. Mỗi buổi học Toán kéo dài $20$ phút, mỗi buổi học Tiếng Anh kéo dài $45$ phút. Vì còn phải tham gia câu lạc bộ thể thao, tổng thời gian tự học hai môn này trong tuần của Nhung không vượt quá $630$ phút. Gọi $x$, $y$ lần lượt là số buổi học Toán và Tiếng Anh trong tuần. Nếu Nhung đã lên lịch $3$ buổi học Toán thì số buổi học Tiếng Anh nhiều nhất có thể là}
  {$12$.}
  {$28$.}{$11$.}{$13$.}

\cauTN[id=T10-C2-B3-D08-VD-074,muc=VD,loigiai={Điều kiện: $15x+45y\le 360$. Với $x=5$: $45y\le 285$, suy ra $y\le \dfrac{19}{3}$. Vậy nhiều nhất $6$ buổi.}]{Trong một tuần, bạn My dành thời gian tự học hai môn Toán và Tiếng Anh. Mỗi buổi học Toán kéo dài $15$ phút, mỗi buổi học Tiếng Anh kéo dài $45$ phút. Vì còn phải tham gia câu lạc bộ thể thao, tổng thời gian tự học hai môn này trong tuần của My không vượt quá $360$ phút. Gọi $x$, $y$ lần lượt là số buổi học Toán và Tiếng Anh trong tuần. Nếu My đã lên lịch $5$ buổi học Toán thì số buổi học Tiếng Anh nhiều nhất có thể là}
  {$6$.}
  {$19$.}{$7$.}{$8$.}

\cauTN[id=T10-C2-B3-D08-VD-075,muc=VD,loigiai={Điều kiện $100x+105y\ge 700$. Với phương án $(0;2)$: $210<700$ nên không đáp ứng.}]{Một bạn học sinh lớp 10 muốn bổ sung năng lượng cho bữa sáng trước giờ học thể dục. Mỗi hộp sữa chua cung cấp khoảng $100$ kcal, mỗi lát bánh mì nguyên cám cung cấp khoảng $105$ kcal. Theo tư vấn của nhân viên y tế trường, bữa sáng của bạn cần cung cấp không ít hơn $700$ kcal. Gọi $x$ là số hộp sữa chua và $y$ là số lát bánh mì bạn ăn. Phương án nào sau đây \textbf{không} đáp ứng yêu cầu về năng lượng?}
  {$0$ hộp sữa chua và $2$ lát bánh mì.}
  {$3$ hộp sữa chua và $4$ lát bánh mì.}{$2$ hộp sữa chua và $5$ lát bánh mì.}{$4$ hộp sữa chua và $3$ lát bánh mì.}

\cauTN[id=T10-C2-B3-D08-VD-076,muc=VD,loigiai={Điều kiện: $25x+15y\le 480$. Với $x=2$: $15y\le 430$, suy ra $y\le \dfrac{86}{3}$. Vậy nhiều nhất $28$ buổi.}]{Trong một tuần, bạn Hà dành thời gian tự học hai môn Toán và Tiếng Anh. Mỗi buổi học Toán kéo dài $25$ phút, mỗi buổi học Tiếng Anh kéo dài $15$ phút. Vì còn phải tham gia câu lạc bộ thể thao, tổng thời gian tự học hai môn này trong tuần của Hà không vượt quá $480$ phút. Gọi $x$, $y$ lần lượt là số buổi học Toán và Tiếng Anh trong tuần. Nếu Hà đã lên lịch $2$ buổi học Toán thì số buổi học Tiếng Anh nhiều nhất có thể là}
  {$28$.}
  {$30$.}{$29$.}{$32$.}

\cauTN[id=T10-C2-B3-D08-VD-077,muc=VD,loigiai={Điều kiện $120x+135y\ge 600$. Với phương án $(1;2)$: $390<600$ nên không đáp ứng.}]{Một bạn học sinh lớp 10 muốn bổ sung năng lượng cho bữa sáng trước giờ học thể dục. Mỗi hộp sữa chua cung cấp khoảng $120$ kcal, mỗi lát bánh mì nguyên cám cung cấp khoảng $135$ kcal. Theo tư vấn của nhân viên y tế trường, bữa sáng của bạn cần cung cấp không ít hơn $600$ kcal. Gọi $x$ là số hộp sữa chua và $y$ là số lát bánh mì bạn ăn. Phương án nào sau đây \textbf{không} đáp ứng yêu cầu về năng lượng?}
  {$1$ hộp sữa chua và $2$ lát bánh mì.}
  {$5$ hộp sữa chua và $1$ lát bánh mì.}{$3$ hộp sữa chua và $3$ lát bánh mì.}{$0$ hộp sữa chua và $5$ lát bánh mì.}

\cauTN[id=T10-C2-B3-D08-VD-078,muc=VD,loigiai={Điều kiện: $25x+15y\le 450$. Với $x=2$: $15y\le 400$, suy ra $y\le \dfrac{80}{3}$. Vậy nhiều nhất $26$ buổi.}]{Trong một tuần, bạn Yến dành thời gian tự học hai môn Toán và Tiếng Anh. Mỗi buổi học Toán kéo dài $25$ phút, mỗi buổi học Tiếng Anh kéo dài $15$ phút. Vì còn phải tham gia câu lạc bộ thể thao, tổng thời gian tự học hai môn này trong tuần của Yến không vượt quá $450$ phút. Gọi $x$, $y$ lần lượt là số buổi học Toán và Tiếng Anh trong tuần. Nếu Yến đã lên lịch $2$ buổi học Toán thì số buổi học Tiếng Anh nhiều nhất có thể là}
  {$26$.}
  {$16$.}{$30$.}{$25$.}

\cauTN[id=T10-C2-B3-D08-VD-079,muc=VD,loigiai={Điều kiện: $20x+40y\le 720$. Với $x=5$: $40y\le 620$, suy ra $y\le \dfrac{31}{2}$. Vậy nhiều nhất $15$ buổi.}]{Trong một tuần, bạn Mai dành thời gian tự học hai môn Toán và Tiếng Anh. Mỗi buổi học Toán kéo dài $20$ phút, mỗi buổi học Tiếng Anh kéo dài $40$ phút. Vì còn phải tham gia câu lạc bộ thể thao, tổng thời gian tự học hai môn này trong tuần của Mai không vượt quá $720$ phút. Gọi $x$, $y$ lần lượt là số buổi học Toán và Tiếng Anh trong tuần. Nếu Mai đã lên lịch $5$ buổi học Toán thì số buổi học Tiếng Anh nhiều nhất có thể là}
  {$15$.}
  {$18$.}{$16$.}{$31$.}

\cauTN[id=T10-C2-B3-D08-VD-080,muc=VD,loigiai={Điều kiện $70x+150y\ge 500$. Với phương án $(4;0)$: $280<500$ nên không đáp ứng.}]{Một bạn học sinh lớp 10 muốn bổ sung năng lượng cho bữa sáng trước giờ học thể dục. Mỗi hộp sữa chua cung cấp khoảng $70$ kcal, mỗi lát bánh mì nguyên cám cung cấp khoảng $150$ kcal. Theo tư vấn của nhân viên y tế trường, bữa sáng của bạn cần cung cấp không ít hơn $500$ kcal. Gọi $x$ là số hộp sữa chua và $y$ là số lát bánh mì bạn ăn. Phương án nào sau đây \textbf{không} đáp ứng yêu cầu về năng lượng?}
  {$4$ hộp sữa chua và $0$ lát bánh mì.}
  {$3$ hộp sữa chua và $3$ lát bánh mì.}{$1$ hộp sữa chua và $5$ lát bánh mì.}{$4$ hộp sữa chua và $2$ lát bánh mì.}

\cauTN[id=T10-C2-B3-D08-VD-081,muc=VD,loigiai={Điều kiện: $15x+25y\le 600$. Với $x=5$: $25y\le 525$, suy ra $y\le 21$. Vậy nhiều nhất $21$ buổi.}]{Trong một tuần, bạn Lan dành thời gian tự học hai môn Toán và Tiếng Anh. Mỗi buổi học Toán kéo dài $15$ phút, mỗi buổi học Tiếng Anh kéo dài $25$ phút. Vì còn phải tham gia câu lạc bộ thể thao, tổng thời gian tự học hai môn này trong tuần của Lan không vượt quá $600$ phút. Gọi $x$, $y$ lần lượt là số buổi học Toán và Tiếng Anh trong tuần. Nếu Lan đã lên lịch $5$ buổi học Toán thì số buổi học Tiếng Anh nhiều nhất có thể là}
  {$21$.}
  {$24$.}{$22$.}{$35$.}

\cauTN[id=T10-C2-B3-D08-VD-082,muc=VD,loigiai={Điều kiện $80x+105y\ge 600$. Với phương án $(2;1)$: $265<600$ nên không đáp ứng.}]{Một bạn học sinh lớp 10 muốn bổ sung năng lượng cho bữa sáng trước giờ học thể dục. Mỗi hộp sữa chua cung cấp khoảng $80$ kcal, mỗi lát bánh mì nguyên cám cung cấp khoảng $105$ kcal. Theo tư vấn của nhân viên y tế trường, bữa sáng của bạn cần cung cấp không ít hơn $600$ kcal. Gọi $x$ là số hộp sữa chua và $y$ là số lát bánh mì bạn ăn. Phương án nào sau đây \textbf{không} đáp ứng yêu cầu về năng lượng?}
  {$2$ hộp sữa chua và $1$ lát bánh mì.}
  {$2$ hộp sữa chua và $5$ lát bánh mì.}{$4$ hộp sữa chua và $3$ lát bánh mì.}{$1$ hộp sữa chua và $5$ lát bánh mì.}

\cauDS[id=T10-C2-B3-D08-VD-083,muc=VD,loigiai={Khối lượng hàng $4x+5y\le 33$. Thay số để kiểm tra từng phương án; chỉ chở xoài: $4x\le 33$, $x\le 8$. a) Đúng. b) Sai. c) Đúng. d) Đúng.}]{Một chiếc xe tải nhỏ của hợp tác xã chở hai loại hàng nông sản: thùng xoài nặng $4$ kg/thùng và thùng thanh long nặng $5$ kg/thùng (đã tính bằng đơn vị chục kg). Tải trọng cho phép của xe không vượt quá $33$ chục kg. Gọi $x$, $y$ lần lượt là số thùng xoài và thùng thanh long xếp lên xe ($x,y\in\mathbb{N}$).}{%
  \y{Điều kiện về tải trọng là $4x+5y\le 33$.}{d}
  \y{Xe có thể chở $6$ thùng xoài và $6$ thùng thanh long.}{s}
  \y{Nếu chỉ chở xoài thì xe chở được tối đa $8$ thùng.}{d}
  \y{Nếu xe chở $4$ thùng xoài thì số thùng thanh long chở thêm được nhiều nhất là $3$.}{d}}

\cauDS[id=T10-C2-B3-D08-VD-084,muc=VD,loigiai={Khối lượng hàng $3x+2y\le 40$. Thay số để kiểm tra từng phương án; chỉ chở xoài: $3x\le 40$, $x\le 13$. a) Đúng. b) Đúng. c) Đúng. d) Đúng.}]{Một chiếc xe tải nhỏ của hợp tác xã chở hai loại hàng nông sản: thùng xoài nặng $3$ kg/thùng và thùng thanh long nặng $2$ kg/thùng (đã tính bằng đơn vị chục kg). Tải trọng cho phép của xe không vượt quá $40$ chục kg. Gọi $x$, $y$ lần lượt là số thùng xoài và thùng thanh long xếp lên xe ($x,y\in\mathbb{N}$).}{%
  \y{Điều kiện về tải trọng là $3x+2y\le 40$.}{d}
  \y{Xe có thể chở $3$ thùng xoài và $3$ thùng thanh long.}{d}
  \y{Nếu chỉ chở xoài thì xe chở được tối đa $13$ thùng.}{d}
  \y{Nếu xe chở $6$ thùng xoài thì số thùng thanh long chở thêm được nhiều nhất là $11$.}{d}}

\cauDS[id=T10-C2-B3-D08-VD-085,muc=VD,loigiai={Khối lượng hàng $2x+3y\le 27$. Thay số để kiểm tra từng phương án; chỉ chở xoài: $2x\le 27$, $x\le 13$. a) Đúng. b) Sai. c) Đúng. d) Đúng.}]{Một chiếc xe tải nhỏ của hợp tác xã chở hai loại hàng nông sản: thùng xoài nặng $2$ kg/thùng và thùng thanh long nặng $3$ kg/thùng (đã tính bằng đơn vị chục kg). Tải trọng cho phép của xe không vượt quá $27$ chục kg. Gọi $x$, $y$ lần lượt là số thùng xoài và thùng thanh long xếp lên xe ($x,y\in\mathbb{N}$).}{%
  \y{Điều kiện về tải trọng là $2x+3y\le 27$.}{d}
  \y{Xe có thể chở $1$ thùng xoài và $9$ thùng thanh long.}{s}
  \y{Nếu chỉ chở xoài thì xe chở được tối đa $13$ thùng.}{d}
  \y{Nếu xe chở $6$ thùng xoài thì số thùng thanh long chở thêm được nhiều nhất là $5$.}{d}}

\cauDS[id=T10-C2-B3-D08-VD-086,muc=VD,loigiai={Khối lượng hàng $5x+3y\le 36$. Thay số để kiểm tra từng phương án; chỉ chở xoài: $5x\le 36$, $x\le 7$. a) Đúng. b) Đúng. c) Đúng. d) Đúng.}]{Một chiếc xe tải nhỏ của hợp tác xã chở hai loại hàng nông sản: thùng xoài nặng $5$ kg/thùng và thùng thanh long nặng $3$ kg/thùng (đã tính bằng đơn vị chục kg). Tải trọng cho phép của xe không vượt quá $36$ chục kg. Gọi $x$, $y$ lần lượt là số thùng xoài và thùng thanh long xếp lên xe ($x,y\in\mathbb{N}$).}{%
  \y{Điều kiện về tải trọng là $5x+3y\le 36$.}{d}
  \y{Xe có thể chở $2$ thùng xoài và $4$ thùng thanh long.}{d}
  \y{Nếu chỉ chở xoài thì xe chở được tối đa $7$ thùng.}{d}
  \y{Nếu xe chở $3$ thùng xoài thì số thùng thanh long chở thêm được nhiều nhất là $7$.}{d}}

\cauTLN[id=T10-C2-B3-D08-VD-087,muc=VD,loigiai={Đếm các cặp số nguyên dương $(x;y)$ thoả mãn $2x+3y\le 18$: có $21$ cặp.}]{Lớp 10A1 chuẩn bị quà cho buổi giao lưu với trường bạn. Mỗi túi quà loại I có giá $2$ chục nghìn đồng, mỗi túi quà loại II có giá $3$ chục nghìn đồng. Quỹ lớp dành cho việc mua quà không vượt quá $18$ chục nghìn đồng và lớp muốn mua \textbf{mỗi loại ít nhất một túi}. Hỏi lớp có bao nhiêu cách chọn số lượng túi quà mỗi loại?}{21}

\cauTLN[id=T10-C2-B3-D08-VD-088,muc=VD,loigiai={Xét lần lượt $x=0,1,\ldots,11$ và đếm số giá trị $y$ tự nhiên tương ứng: 3+3+3+3+2+2+2+2+1+1+1+1=24.}]{Có bao nhiêu cặp số tự nhiên $(x;y)$ thoả mãn bất phương trình $x+4y\le 11$?}{24}

\cauTLN[id=T10-C2-B3-D08-VD-089,muc=VD,loigiai={Đếm các cặp số nguyên dương $(x;y)$ thoả mãn $4x+5y\le 22$: có $8$ cặp.}]{Lớp 10A5 chuẩn bị quà cho buổi giao lưu với trường bạn. Mỗi túi quà loại I có giá $4$ chục nghìn đồng, mỗi túi quà loại II có giá $5$ chục nghìn đồng. Quỹ lớp dành cho việc mua quà không vượt quá $22$ chục nghìn đồng và lớp muốn mua \textbf{mỗi loại ít nhất một túi}. Hỏi lớp có bao nhiêu cách chọn số lượng túi quà mỗi loại?}{8}

\cauTLN[id=T10-C2-B3-D08-VD-090,muc=VD,loigiai={Xét lần lượt $x=0,1,\ldots,1$ và đếm số giá trị $y$ tự nhiên tương ứng: 4+2=6.}]{Có bao nhiêu cặp số tự nhiên $(x;y)$ thoả mãn bất phương trình $4x+2y\le 7$?}{6}

\cauTN[id=T10-C2-B3-D08-VDC-091,muc=VDC,loigiai={Với mỗi $y$, đếm các $x$ tự nhiên thoả mãn $y\le x\le \dfrac{25-2y}{5}$; cộng lại được $14$ cặp.}]{Có bao nhiêu cặp số tự nhiên $(x;y)$ thoả mãn đồng thời $5x+2y\le 25$ và $x\ge y$?}
  {$14$.}
  {$17$.}{$16$.}{$11$.}

\cauTN[id=T10-C2-B3-D08-VDC-092,muc=VDC,loigiai={Điều kiện $7x+3y\le 33$, $x,y\ge1$. Để $x+y$ lớn nhất, ưu tiên loại rẻ hơn; kiểm tra các trường hợp được tối đa $9$ bộ (ví dụ $x=1$, $y=8$).}]{Thư viện trường dành không quá $33$ triệu đồng để mua sách tham khảo môn Toán (giá $7$ triệu đồng/bộ) và sách tham khảo môn Văn (giá $3$ triệu đồng/bộ). Thư viện muốn mua mỗi loại ít nhất một bộ và tổng số bộ sách là lớn nhất. Tổng số bộ sách lớn nhất mua được là}
  {$9$.}
  {$10$.}{$6$.}{$8$.}

\cauTN[id=T10-C2-B3-D08-VDC-093,muc=VDC,loigiai={Với mỗi $y$, đếm các $x$ tự nhiên thoả mãn $y\le x\le \dfrac{19-3y}{4}$; cộng lại được $11$ cặp.}]{Có bao nhiêu cặp số tự nhiên $(x;y)$ thoả mãn đồng thời $4x+3y\le 19$ và $x\ge y$?}
  {$11$.}
  {$15$.}{$12$.}{$8$.}

\cauTN[id=T10-C2-B3-D08-VDC-094,muc=VDC,loigiai={Điều kiện $6x+3y\le 33$, $x,y\ge1$. Để $x+y$ lớn nhất, ưu tiên loại rẻ hơn; kiểm tra các trường hợp được tối đa $10$ bộ (ví dụ $x=1$, $y=9$).}]{Thư viện trường dành không quá $33$ triệu đồng để mua sách tham khảo môn Toán (giá $6$ triệu đồng/bộ) và sách tham khảo môn Văn (giá $3$ triệu đồng/bộ). Thư viện muốn mua mỗi loại ít nhất một bộ và tổng số bộ sách là lớn nhất. Tổng số bộ sách lớn nhất mua được là}
  {$10$.}
  {$8$.}{$13$.}{$9$.}

\cauTN[id=T10-C2-B3-D08-VDC-095,muc=VDC,loigiai={Với mỗi $y$, đếm các $x$ tự nhiên thoả mãn $y\le x\le \dfrac{19-3y}{2}$; cộng lại được $26$ cặp.}]{Có bao nhiêu cặp số tự nhiên $(x;y)$ thoả mãn đồng thời $2x+3y\le 19$ và $x\ge y$?}
  {$26$.}
  {$28$.}{$29$.}{$25$.}

\cauTN[id=T10-C2-B3-D08-VDC-096,muc=VDC,loigiai={Điều kiện $7x+5y\le 38$, $x,y\ge1$. Để $x+y$ lớn nhất, ưu tiên loại rẻ hơn; kiểm tra các trường hợp được tối đa $7$ bộ (ví dụ $x=1$, $y=6$).}]{Thư viện trường dành không quá $38$ triệu đồng để mua sách tham khảo môn Toán (giá $7$ triệu đồng/bộ) và sách tham khảo môn Văn (giá $5$ triệu đồng/bộ). Thư viện muốn mua mỗi loại ít nhất một bộ và tổng số bộ sách là lớn nhất. Tổng số bộ sách lớn nhất mua được là}
  {$7$.}
  {$5$.}{$9$.}{$4$.}

\cauDS[id=T10-C2-B3-D08-VDC-097,muc=VDC,loigiai={Liệt kê theo $x$: số nghiệm là $30$; với $y=0$ có $x=0,\ldots,11$; giá trị lớn nhất của $x+2y$ là $11$. a) Đúng. b) Đúng. c) Đúng. d) Sai.}]{Xét các nghiệm là cặp số tự nhiên của bất phương trình $x+3y\le 11$.}{%
  \y{Có đúng $30$ cặp số tự nhiên $(x;y)$ là nghiệm.}{d}
  \y{Có $12$ nghiệm tự nhiên với $y=0$.}{d}
  \y{Giá trị lớn nhất của $x+2y$ với $(x;y)$ là nghiệm tự nhiên bằng $11$.}{d}
  \y{Cặp số $(11;1)$ là nghiệm của bất phương trình.}{s}}

\cauDS[id=T10-C2-B3-D08-VDC-098,muc=VDC,loigiai={Liệt kê theo $x$: số nghiệm là $30$; với $y=0$ có $x=0,\ldots,4$; giá trị lớn nhất của $x+2y$ là $18$. a) Đúng. b) Đúng. c) Đúng. d) Đúng.}]{Xét các nghiệm là cặp số tự nhiên của bất phương trình $2x+y\le 9$.}{%
  \y{Có đúng $30$ cặp số tự nhiên $(x;y)$ là nghiệm.}{d}
  \y{Có $5$ nghiệm tự nhiên với $y=0$.}{d}
  \y{Giá trị lớn nhất của $x+2y$ với $(x;y)$ là nghiệm tự nhiên bằng $18$.}{d}
  \y{Cặp số $(4;1)$ là nghiệm của bất phương trình.}{d}}

\cauTLN[id=T10-C2-B3-D08-VDC-099,muc=VDC,loigiai={Liệt kê theo từng giá trị của $x$ và chỉ giữ các cặp có $x+y$ chẵn; được $9$ cặp.}]{Có bao nhiêu cặp số tự nhiên $(x;y)$ thoả mãn $3x+3y\le 15$ và $x+y$ là số chẵn?}{9}

\cauTLN[id=T10-C2-B3-D08-VDC-100,muc=VDC,loigiai={Xét lần lượt $x=0,1,\ldots,6$ và đếm số giá trị $y$ tự nhiên tương ứng: 2+2+2+1+1+1+1=10.}]{Có bao nhiêu cặp số tự nhiên $(x;y)$ thoả mãn bất phương trình $x+4y\le 6$?}{10}
